\documentclass[10pt,a4paper]{article}
\usepackage[utf8]{inputenc}
\usepackage[ngerman]{babel}
\usepackage{amsmath}
\usepackage{amsfonts}
\usepackage{amssymb}
\usepackage{mathtools}
\usepackage{microtype}
\usepackage{authblk}
\usepackage{titlesec}

\usepackage{xcolor}
\usepackage{empheq}
\usepackage{tikz-cd} 

% section title formating
\renewcommand{\thesection}{\Roman{section}}
\titleformat{\section}[block]{\normalsize\bfseries}{\thesection.}{1em}{\MakeUppercase}
\titlespacing*{\section}{0pt}{11pt plus 2pt minus 2pt}{11pt plus 2pt minus 2pt}

\renewcommand{\thesubsection}{\arabic{section}.\Alph{subsection}}
\titleformat{\subsection}{\normalsize\bfseries}{\thesubsection}{1em}{}
\titlespacing*{\subsection}{0pt}{11pt plus 2pt minus 2pt}{11pt plus 2pt minus 2pt}

\renewcommand{\thesubsubsection}{\thesubsection.\arabic{subsubsection}}
\titleformat{\subsubsection}{\normalsize\itshape}{\thesubsubsection}{1em}{}
\titlespacing*{\subsubsection}{0pt}{11pt plus 2pt minus 2pt}{11pt plus 2pt minus 2pt}
   
% custom macros
\newcommand{\mydef}[1]{\textls{#1}}
\newcommand{\nul}{\textrm{---}}
\newcommand{\trans}{^\mathsf{T}}

\newcommand{\ii}{\mathrm{i}}
\newcommand{\dt}{\mathrm{d}}
\newcommand{\Dt}{\mathrm{D}}
\newcommand{\Tt}{\mathrm{T}}
\newcommand{\ad}{\mathrm{ad}}
\newcommand{\Ad}{\mathrm{Ad}}

\newcommand{\cZ}{\mathcal{Z}}
\newcommand{\cS}{\mathcal{S}}

\newcommand*\resultbox[1]{%
  \colorbox{blue!10}{\hspace{1em}#1\hspace{1em}}}
\newcommand*\definitionbox[1]{%
  \colorbox{yellow!20}{\hspace{1em}#1\hspace{1em}}}


% %%%%%%%
% title %
% %%%%%%%

\title{Zu Analysis und Zusammenhängen von Lie-Gruppen}
\author{Hendrik Meier}

\affil{Oberuhldingen}

\affil{}
\renewcommand\Affilfont{\itshape}

\date{\today}


% %%%%%%%%%%
% document %
% %%%%%%%%%%

\begin{document}
\maketitle

Dieser Aufsatz diskutiert auf einführender Weise Aspekte der Analysis und Zusammenhangstheorie auf Lie-Gruppen.
Motivation ist dabei, im Kontext der Navigation zu klären, welcher Zusammenhang den Transport von Fehlervektoren bzw. Kovarianzen der Lage geometrisch adäquat beschreibt.

Es ist dem Autoren dabei ein Anliegen, die Notation möglichst transparent zu halten --- zum Beispiel in Bezug auf die Verwendung der Klammer $[-,-]$,
die in der differenzialgeometrischen Literatur etwa sowohl für Ableitungen als auch für die algebraische Lie-Klammer verwendet wird.

\section{Differenzialrechnung auf Lie-Gruppen}

\subsection{Parallelisierung von Lie-Gruppen}

Sei $G$ ein Lie-Gruppe der Dimension~$n$.\footnote{
	In diesem Aufsatz betrachten wir ausschließlich glatte ($C^\infty$) Lie-Gruppen.
}
Mithilfe der \emph{Rechtsmultiplikation}
parallelisieren wir $G$, indem wir an einem Punkt $g\in G$ angeheftete Tangentialvektoren $\xi_g\in\Tt_gG$ mit einem Vektor im Tangentialraum $\xi=\xi_e\in\Tt_eG$ des neutralen Elements~$e$ identifizieren,
\begin{align}
	\xi = \Tt R_g|_e^{-1}\cdot\xi_g
	\ ,
\end{align}
wobei
\begin{align}
	R_g: G &\rightarrow G
	\ ,\nonumber\\
	h &\mapsto hg
\end{align}
die Rechtsmultiplikation und $\Tt R_g|_h: \Tt_hG\rightarrow\Tt_{hg}G$ ihre Ableitung bezeichnen.\footnote{
Für (koordinatenfrei geschriebene) Ableitungen auf Mannigfaltigkeiten verwenden wir, dem Lehrbuch \cite{StorchWiebe2005} folgend, das Symbol $\Tt$.
Das Symbol $\Dt$ beschränken wir auf Ableitungen von Abbildungen zwischen Vektorräumen, wozu insbesondere auch $\mathfrak{g}=\Tt_eG$ gehört.
}

Im Folgenden nennen wir den Tangentialraum $\Tt_eG$ am neutralen Element
\begin{empheq}[box=\definitionbox]{align}
	\mathfrak{g} := \Tt_eG
	\ .
\end{empheq}
Das praktische Rechnen mit Tangentialvektoren wird ausschließlich in $\mathfrak{g}$ erfolgen,
wobei wir jedoch in Formeln und Herleitungen transparent zwischen $\mathfrak{g}$ und einem durch $\mathfrak{g}$ repräsentierten Tangentialraum $\Tt_gG$ unterscheiden wollen.

Sei umgekehrt $\xi\in\Tt_eG$ ein am neutralen Element angehefteter Tangentialvektor.
Über die Rechtsmultiplikation definieren wir das \mydef{rechtsinvariante} Vektorfeld
\begin{align}
	X_\xi: G &\rightarrow \Tt G
	\ ,\nonumber\\
	g &\mapsto \Tt R_g|_e\cdot\xi \in\Tt_g G
	\ .
\end{align}
Ein Vektorfeld $X$ nennen wir genau dann \mydef{rechtsinvariant}, wenn
\begin{align}
	X|_{hg} &= \Tt R_g|_h\cdot X|_{h}
\end{align}
für alle $g\in G$ gilt.
Ein rechtsinvariantes Vektorfeld~$X$ ist über seinen Wert $X|_e$ am neutralen Element offenbar eindeutig bestimmt und kann daher mit dem Tangentialvektor $\xi=X|_e\in\Tt_eG$ identifiziert werden.

\subsection{Differenziation auf Lie-Gruppen}

Die Ableitung einer differenzierbaren Abbildung $F:G\rightarrow G$ an einer Stelle $g\in G$ ist zunächst die Abbildung 
\begin{align*}
	\Tt F|_g:
	\Tt_gG\rightarrow\Tt_{F(g)}G
	\ .
\end{align*}
Zum praktischen Rechnen identifizieren wir mit Hilfe der Parallelisierung Urbild- und Bildraum mit $\mathfrak{g}=\Tt_eG$.
Die auf diese Weise in $\mathfrak{g}$ behandelte Ableitung bezeichnen wir mit ``$\Dt$'',
\begin{empheq}[box=\definitionbox]{align}
	\Dt F|_g: 
	\mathfrak{g} &\rightarrow \mathfrak{g}
	\ ,\nonumber\\
	\xi &\mapsto 
	\Tt R_{F(g)}^{-1}|_{F(g)}
	\cdot \Tt F|_g
	\cdot \Tt R_{g}|_e
	\ ,
\end{empheq}
wobei das letzte Verkettungsglied auch als $\Tt R_{F(g)}^{-1}|_{F(g)}=\Tt R_{F(g)}|_e^{-1}$ geschrieben werden kann.

\subsubsection{Multiplikation}

Es bezeichne 
\begin{align}
	\mathrm{mult}:G\times G\rightarrow G
\end{align}
die Multiplikationsabbildung $(g,h)\mapsto gh$,
wobei das Produkt $G\times G$ seine topologische, analytische und algebraische Struktur von $G$ erbt.
Zur Bestimmung von $\Dt\,\mathrm{mult}|_{g,h}$ seien für $\varepsilon>0$ die Kurvenkeime $\gamma,\eta: (-\varepsilon,\varepsilon)\rightarrow G$ mit $\gamma(0)=g$ und $\dot{\gamma}(0)=\xi_g\in\Tt_gG$ sowie $\eta(0)=h$ und $\dot{\eta}(0)=\upsilon_h\in\Tt_hG$ gegeben.

Damit berechnen wir einerseits
\begin{align*}
	\Tt\,\mathrm{mult}|_{(g,h)}\cdot (\xi_g, 0)
	&= \frac{\dt}{\dt t}
		R_h\gamma(t)
	\Big|_{t=0}
	= \Tt R_h|_g\cdot\xi_g
	\in \Tt_{gh}G
	\ ,
\end{align*}
andererseits
\begin{align*}
	\Tt\,\mathrm{mult}|_{(g,h)}\cdot (0, \upsilon_h)
	&= \frac{\dt}{\dt t}
		L_g\eta(t)
	\Big|_{t=0}
	= \Tt L_g|_h\cdot\upsilon_h
	\in \Tt_{gh}G
	\ ,
\end{align*}
wobei $L_g:h\mapsto gh$ die Linksmultiplikation bezeichnet.
Wir finden mit
\begin{empheq}[box=\resultbox]{align}
	\Tt\,\mathrm{mult}|_{(g,h)}\cdot (\xi_g, \upsilon_h)
	&= \Tt R_h|_g\cdot\xi_g + \Tt L_g|_h\cdot\upsilon_h
	\ .
\end{empheq}
die Produktregel für 
$\Tt\,\mathrm{mult}|_{(g,h)}:\Tt_gG\times\Tt_hG\rightarrow\Tt_{hg}G$.

Aufgefasst als Abbildung $\mathfrak{g}\times\mathfrak{g}\rightarrow\mathfrak{g}$ finden wir
\begin{align*}
	\Dt\,\mathrm{mult}|_{(g,h)}\cdot (\xi, 0)
	&= 
	\Tt R_{gh}^{-1}|_{gh}
	\cdot
	\Tt R_h|_g
	\cdot
	\Tt R_g|_g
	\cdot\xi
	\nonumber\\
	&= \Tt (R_{(gh)^{-1}}\circ R_h\circ R_g)|_e\cdot\xi
	= \xi
\end{align*}
und
\begin{align*}
	\Dt\,\mathrm{mult}|_{(g,h)}\cdot (0, \upsilon)
	&=
	\Tt R_{gh}^{-1}|_{gh}
	\cdot	
	\Tt L_g|_h
	\cdot
	\Tt R_h|_e
	\cdot
	\upsilon
	=
	\Tt
	(R_{g^{-1}}\circ L_g)|_e
	\cdot
	\upsilon	
	\ .
\end{align*}
Zusammengefasst ist also
\begin{empheq}[box=\resultbox]{align}
	\Dt\,\mathrm{mult}|_{(g, h)} \cdot (\xi, \upsilon)
	&=
	\xi + \Ad(g)\cdot\upsilon
	\ ,
\end{empheq}
wobei
\begin{empheq}[box=\definitionbox]{align}
\label{eq:Ad}
	\Ad: G&\rightarrow \mathrm{GL}(\mathfrak{g})
	\ ,\nonumber\\
	g &\mapsto
	\Tt
	(R_{g^{-1}}L_g)|_e
\end{empheq}
die \mydef{adjungierte Darstellung} von $G$ bezeichnet.

\subsubsection{Inversenbildung}

Sei
\begin{align}
	\mathrm{inv}:G\rightarrow G
\end{align} 
die Abbildung der Inversenbildung $g\mapsto g^{-1}$.

Wir drücken dazu die \emph{konstante} Abbildung auf das neutrale Element als Komposition
\begin{alignat}{2}
	G 
	&\rightarrow G\times G 
	&&\rightarrow G
	\ ,\nonumber\\
	g 
	&\mapsto
	\begin{pmatrix}
		\mathrm{inv}(g) \\ g
	\end{pmatrix}
	&&\mapsto \mathrm{mult}(\mathrm{inv}(g), g) = e
\end{alignat}
aus.
Für $g\in G$ in $\xi_g\in\Tt_gG$ gilt dann nach der Kettenregel
\begin{align*}
	\mathfrak{g}\ni	
	0 &= 
	\Tt\,\mathrm{mult}|_{(g^{-1},g)}
	\cdot
	\begin{pmatrix}
		\Tt\,\mathrm{inv}|_{g} \cdot \xi_g \\
		\xi_g
	\end{pmatrix}
	\nonumber\\
	&= 
	\Tt R_g|_{g^{-1}}\cdot \Tt\,\mathrm{inv}|_{g} \cdot \xi_g
	+ \Tt L_{g^{-1}}|_g\cdot \xi_g
	\ .
\end{align*}
Folglich finden wir mit $\xi_g=\Tt R_g|_e\cdot \xi$ für die Ableitung $\Dt\,\mathrm{inv}|_{g}:\mathfrak{g}\rightarrow\mathfrak{g}$
\begin{align*}
	\Dt\,\mathrm{inv}|_{g} \cdot \xi
	&= -\Tt L_{g^{-1}}|_g\cdot \Tt R_g|_e \cdot \xi
	= -\Tt (L_{g^{-1}}\circ R_g)|_e \cdot \xi
	\ ,
\end{align*}
also
\begin{empheq}[box=\resultbox]{align}
	\Dt\,\mathrm{inv}|_{g}
	&= -\Ad(g^{-1})
	\ .
\end{empheq}

\subsection{Vektorfelder auf Mannigfaltigkeiten}

Wir tragen in diesem Abschnitt grundlegende Eigenschaften von Vektorfeldern auf Mannigfaltigkeiten zusammen.\footnote{
	Wir betrachten durchweg glatte ($C^\infty$) Mannigfaltigkeiten.
}
Sei dazu $X$ allgemein ein Vektorfeld einer Mannigfaltigkeit $M$ der Dimension~$n$.

\subsubsection{Push-Forward eines Vektorfelds}

Sei $F:M\rightarrow N$ ein Diffeomorphismus auf die Mannigfaltigkeit~$N$.
Mit $F$ können wir $X$ als \mydef{Push-Forward}  auf $N$ transportieren,
\begin{align}
	F_*X: N &\rightarrow \Tt N
	\ ,\nonumber\\
	q &\mapsto \Tt F|_{F^{-1}(q)}\cdot X|_{F^{-1}(q)}
	\ .
\end{align}

\subsubsection{Fluss eines Vektorfelds}

Integralkurven $t\mapsto\gamma(t)\in M$ zu einem Vektorfeld~$X$ sind Kurven, deren Geschwindigkeiten durch Vektorfeld gegeben sind,
\begin{align*}
	\dot{\gamma}(t) &= X|_{\gamma(t)}
	\ .
\end{align*}
Flüsse bieten eine übersichtliche Darstellung über Integralkurven.

Aus der Theorie gewöhnlicher Differenzialgleichungen folgt die Existenz eines offenen Intervalls $I\subset\mathbb{R}$ mit $0\in I$ und für jedes $p_0\in M$ die Existenz einer offenen Menge $V\subset M$ mit $p_0\in V$, auf deren Produkt der \mydef{Fluss}
\begin{align}
	\Phi_X: I\times V \rightarrow M
\end{align}
mit
\begin{align}
	\frac{\dt}{\dt t}\Phi_X(t; p) = X|_{\Phi_X(t; p)}
\end{align}
wohldefiniert und glatt ist.
Die Kurve $t\mapsto\Phi_X(t; p)$ ist für $p$ fix die eindeutige Integralkurve durch $p$.

Der Fluss $\Phi_X$ des (zeitunabhängigen!) Vektorfelds~$X$ erfüllt offenbar die \mydef{Halbgruppeneigenschaft} im Sinne von
\begin{align}
	\Phi_X(t_1 + t_2; -)
	&= \Phi_X(t_2; \Phi_X(t_1; -))
	\ ,
\end{align}
sofern beide Seiten existieren.
Außerdem ist $\Phi_X(0; -)=\mathrm{id}$.


\subsubsection{Lie-Produkt von Vektorfeldern}

Für $t\in I$ ist $\Phi_X(t):=\Phi_X(t;-)$ ein (mindestens) \emph{lokaler} Diffeomorphismus auf sein Bild.
Seine Umkehrung bezeichnen wir mit $\Phi_X^{-1}(t):q\mapsto\Phi_X^{-1}(t;q)$, wobei $\Phi_X^{-1}(t;q)=\Phi_X(-t;q)$.
Sei nun $Y$ ein weiteres Vektorfeld auf $M$.

Das \mydef{Lie-Produkt} der Vektorfelder $X$ und $Y$ definieren wir als
\begin{align}
	[X, Y]\big|_{p} &:= 
	\frac{\dt}{\dt t}
	\big(
		[\Phi_X(t)]_* Y
	\big)
	\big|_{p}
	\Big|_{t=0}
\end{align}
für alle $p\in M$.

Ist beispielsweise $M=\mathbb{R}^n$, was \emph{lokal} durch Wahl einer Karte stets erreicht werden kann, finden wir für
\begin{align*}
	[X, Y]\big|_{p}
	&=
	\frac{\dt}{\dt t}
	\Big(
		\Dt\Phi_X(t)|_{\Phi_X^{-1}(t, p)}
		\cdot
		Y|_{\Phi_X^{-1}(t, p)}
	\Big)\Big|_{t=0}
\end{align*}
unter Verwendung von
\begin{align*}
	\Phi_X(t; p)
	&= p + t X|_{p} + \mathrm{O}(t^2)
	\ ,\nonumber\\
	\Phi_X^{-1}(t; p)
	&= p - t X|_{p} + \mathrm{O}(t^2)
	\ ,\nonumber\\
	\Dt\Phi_X(t)\big|_{p}
	&= \mathrm{id} + t\,\Dt X|_{p} + \mathrm{O}(t^2)
	\ ,\nonumber\\
	Y|_{\Phi_X^{-1}(t, p)}
	&= Y|_{p} - t\,\Dt Y|_{p}\cdot X|_{p} + \mathrm{O}(t^2)
\end{align*}
explizit
\begin{align}
	[X, Y]\big|_{p}
	&=
	\Dt X|_{p}
	\cdot
	Y|_{p}
	-
	\Dt Y|_{p}
	\cdot
	X|_{p}
	\ .
\end{align}
Das Lie-Produkt entspricht damit der \mydef{Lie-Klammer} von Vektorfeldern.
Die Jacobi-Identität der Lie-Klammer, die dem Raum der Vektorfelder die algebraische Struktur einer Lie-Algebra gibt, bestätigt man durch direktes Überprüfen.

Die Lie-Klammer ist auch im Mannigfaltigkeitsfall unabhängig von der Wahl einer Karte:
Ist $F$ ein Diffeomorphismus auf $M=\mathbb{R}^n$, so ist nämlich
\begin{align}
	[F_*X, F_*Y]\big|_{p}
	&= 
	\Dt \big(
		\Dt F|_{F^{-1}(p)} \cdot X|_{F^{-1}(p)}
	\big)\big|_p
	\cdot
	\Dt F|_{F^{-1}(p)} \cdot Y|_{F^{-1}(p)}
	- [X\leftrightarrow Y]
	\nonumber\\
	&=
	\Dt^2F|_{F^{-1}(p)}
	\cdot
	(X|_{F^{-1}(p)}\otimes Y|_{F^{-1}(p)})
	\nonumber\\
	&\quad 
	+
	\Dt F|_{F^{-1}(p)} 
	\cdot \Dt X|_{F^{-1}(p)}
	\cdot Y|_{F^{-1}(p)}
	- [X\leftrightarrow Y]
	\nonumber\\
	&= \Dt F|_{F^{-1}(p)}
	\cdot
	[X, Y]|_{F^{-1}(p)}
	\nonumber\\
	&= F_*[X, Y]|_p
	\ .
\end{align}
Dabei ist verwendet worden, dass für Vektoren $u,v\in\mathbb{R}^n$
\begin{align*}
	\Dt_v(\Dt_u F|_{F^{-1}(p)})\big|_p
	&=
	\Dt_{\Dt F^{-1}|_p\cdot v}
	\Dt_u F\big|_{F^{-1}(p)}
	=:
	\Dt^2F|_{F^{-1}(p)}\cdot(u \otimes \Dt F|_{F^{-1}(p)}^{-1}\cdot v)
\end{align*}
mit $\Dt F^{-1}|_p=\Dt F|_{F^{-1}(p)}^{-1}$ gilt, dass
\begin{align*}
	\Dt \big(
		X|_{F^{-1}(p)}
	\big)\big|_p
	&=
	\Dt X|_{F^{-1}(p)}
	\cdot \Dt F|_{p}^{-1}
\end{align*}
ist sowie dass $\Dt^2$ symmetrisch ist.

\subsection{Lie-Algebren}

\subsubsection{Lie-Klammer auf $\mathfrak{g}$}

Sei wieder $G$ eine Lie-Gruppe der Dimension $n$ und seien $\xi_1,\xi_2\in\mathfrak{g}=\Tt_eG$ zwei Tangentialvektoren am neutralen Element mit zugehörigen rechtsinvarianten Vektorfeldern~$X_{\xi_1}$ und $X_{\xi_2}$.
Die Lie-Klammer $[X_{\xi_1}, X_{\xi_2}]$ ist wegen
\begin{align*}
	(R_g)_*[(R_g)_*X_{\xi_1}, (R_g)_*{\xi_2}]
\end{align*}
ebenfalls rechtsinvariant und identifizierbar mit einem Tangentialvektor in $\Tt_eG$.
Damit trägt $\mathfrak{g}=\Tt_eG$ die Struktur einer \mydef{Lie-Algebra}.
Um die algebraische Lie-Klammer auf $\mathfrak{g}$ von der differenziellen Lie-Klammer $[-,-]$ zwischen beliebigen Vektorfeldern auf $G$ zu unterscheiden, definieren wir für $\xi_1,\xi_2\in\Tt_eG$
\begin{empheq}[box=\definitionbox]{align}
	[-,-]_\mathfrak{g}:
	\mathfrak{g}\times\mathfrak{g}
	&\rightarrow \mathfrak{g}
	\ ,\nonumber\\
	(\xi_1,\xi_2) &\mapsto
	[\xi_1,\xi_2]_\mathfrak{g}
	:= [X_{\xi_1}, X_{\xi_2}]
	\ .
\end{empheq}
$[-,-]_\mathfrak{g}$ ist eine bilineare Abbildung auf $\mathfrak{g}$ mit Werten in $\mathfrak{g}$.

Sei speziell $\xi_1,\ldots,\xi_n$ eine Basis von $\mathfrak{g}$.
Dann lässt sich schreiben
\begin{align}
	[\xi_i,\xi_j]_\mathfrak{g}
	&= \sum_{k=1}^n f^k_{ij} \xi_k
	\ .
\end{align}
Die \mydef{Strukturkonstanten} $f^k_{ij}\in\mathbb{R}$ bezüglich einer gegebenen Basis bestimmen eindeutig die Lie-Algebrastruktur von $\mathfrak{g}$.
	
\subsubsection{Fluss eines rechtsinvarianten Vektorfeldes}

Sei $\xi\in\mathfrak{g}$ und $X_\xi:g\mapsto\Tt R_g|_e\cdot\xi$ das zugehörige rechtsinvariante Vektorfeld.

Sei $\gamma:(-\varepsilon,\varepsilon)\rightarrow G$ eine das neutrale Element passierende lokale Integralkurve zu $X_\xi$,
\begin{align*}
	\gamma(t) &= e
	\\
	\dot{\gamma}(t) &= X_\xi|_{\gamma(t)}
	\ .
\end{align*}
Die Rechnung
\begin{align*}
	\frac{\dt}{\dt t}
	\big(
		L_{\gamma(t)}\,g
	\big)
	\Big|_{t=0}
	&=
	\frac{\dt}{\dt t}
	\big(
		\gamma(t)\,g
	\big)
	\Big|_{t=0}	
	=
	\Tt R_g|_{e} \cdot X_\xi|_{e}
	= X_\xi|_{g}		
\end{align*}
zeigt, dass der Fluss zu $X_\xi$ lokal durch Linksmultiplikation einer Integralkurve durch das neutrale Element gegeben ist,
\begin{align}
	\Phi_{X_\xi}(t; -)
	&= L_{\gamma(t)}
	\ .
\end{align}

\subsubsection{Adjungierte Darstellung}

Die adjungierte Darstellung $\Ad:G\rightarrow\mathrm{GL}(\mathfrak{g})$ nach Gl.~(\ref{eq:Ad}) ist ein Homomorphismus von Lie-Gruppen.
Hier soll seine Ableitung
\begin{empheq}[box=\definitionbox]{align}
	\ad:=\Dt\,\Ad|_e: \mathfrak{g}\rightarrow\mathfrak{gl}(\mathfrak{g})
\end{empheq}
im Tangentialraum $\mathfrak{g}$ des neutralen Elements bestimmt werden.

Nach Definition ist für $g\in G$ und $\upsilon\in\mathfrak{g}$ mit zugehörigem rechtsinvarianten Vektorfeld~$Y_\upsilon$
\begin{align*}
	\Ad(g)\cdot\upsilon
	&= 
	\Tt(L_{g}R_{g^{-1}})|_e\,\upsilon
	=
	\Tt L_{g}|_{g^{-1}}
	\cdot
	\upsilon_{g^{-1}}
	=  
	\big(
		[L_{g}]_* Y_\upsilon
	\big)
	\big|_e
	\ .	
\end{align*}
Sei  $\gamma:(-\varepsilon,\varepsilon)\rightarrow G$ mit $\gamma(0)=e$ und $\dot{\gamma}=X_\xi|_{\gamma(t)}$, insbesondere mit $\dot{\gamma}(0)=\xi\in\mathfrak{g}$	eine Integralkurve zum rechtsinvarianten Vektorfeld $X_\xi$.
Es folgt
\begin{align*}
	\Dt_{\xi}\Ad|_e\cdot\upsilon
	&=
	\frac{\dt}{\dt t}
	\Ad|_{\gamma(t)}\cdot\upsilon
	\Big|_{t=0}
	=
	\frac{\dt}{\dt t}
	\big(
		[L_{\gamma(t)}]_* Y_\upsilon
	\big)
	\big|_e
	\Big|_{t=0}
	\ .
\end{align*}
Da $L_{\gamma(t)}=\Phi_{X_\xi}(t;-)$, folgt
\begin{align*}
	\Dt_{\xi}\Ad|_e\cdot\upsilon
	&=
	[X_\xi,Y_\upsilon] = [\xi, \upsilon]_\mathfrak{g}
	\ .
\end{align*}
$\ad(\xi)$ ist also der durch
\begin{empheq}[box=\resultbox]{align}
	\ad(\xi): \upsilon\mapsto[\xi,\upsilon]_\mathfrak{g}
\end{empheq}
gegebene Endomorphismus auf~$\mathfrak{g}$.

Mit der Jacobi-Identität auf $\mathfrak{g}$ folgt durch
\begin{align*}
	\ad([\xi_1,\xi_2])
	&= [[\xi_1,\xi_2], -]
	= -[[\xi_2,-], \xi_1]-[[-,\xi_1], \xi_2]
	\nonumber\\
	&= \ad(\xi_1)\circ \ad(\xi_2) - \ad(\xi_2)\circ \ad(\xi_1)
	\ ,
\end{align*}
dass $\ad$ ein Lie-Algebrahomomorphismus $\mathfrak{g}\rightarrow\mathfrak{gl}(\mathfrak{g})$ ist.
	
\subsection{Liesche Exponentialabbildung}

\subsubsection{Vorbetrachtungen}

Sei $\xi\in\mathfrak{g}=\Tt_eG$ ein Element der Lie-Algebra zur Lie-Gruppe~$G$ und $X_\xi$ das zugehörige rechtsinvariante Vektorfeld.
Der Fluss $\Phi_{X_\xi}(t;e)$ ist mindestens lokal für ein hinreichend kleines $t=\Delta t$ definiert.
Wir setzen nun $g=\Phi_{X_\xi}(\Delta t;e)$ und 
\begin{align*}
	\tilde{\Phi}_{X_\xi}(t;g)
	&= R_g\circ\Phi_{X_\xi}(t;e)
\end{align*}
Eine einfache Prüfung ergibt, dass $\tilde{\Phi}_{X_\xi}(0;g)=g$ und
\begin{align*}
	\frac{\dt}{\dt t}
		\tilde{\Phi}_{X_\xi}(t;g)
	\Big|_{t=0}
	= \Tt R_g|_e \cdot X_\xi|_e
	= \Tt R_g|_e \cdot \xi
	\ ,
\end{align*}
woraus folgt, dass $\tilde{\Phi}_{X_\xi}(t;g)=\Phi_{X_\xi}(t;g)$.
Der Fluss mit Startpunkt $g$ wird also auf den am neutralen Element~$e$ startenden Fluss zurückgeführt.
Damit impliziert nun die Existenz von $\Phi_{X_\xi}(\Delta t;e)$ jene von $\Phi_{X_\xi}(\Delta t;g)=\Phi_{X_\xi}(2\Delta t;e)$,
woraus wir folgern, dass der Fluss $\Phi_{X_\xi}$ zum rechtsinvarianten Vektorfeld~$X_\xi$ auf ganz $\mathbb{R}\times G$ existiert.

\subsubsection{Definition der Exponentialabbildung}

Die angestellten Vorbetrachtungen begründen die Definition der \mydef{Exponentialabbildung} gemäß
\begin{empheq}[box=\definitionbox]{align}
	\exp: 
	\mathfrak{g} &\rightarrow G
	\ ,\nonumber\\
	\xi &\mapsto \exp\xi := \Phi_{X_\xi}(1; e)
	\ .
\end{empheq}
Umgekehrt ist mit Hilfe der Exponentialabbildung der Fluss eines \emph{rechtsinvarianten} Vektorfelds $X_\xi$ mit $\xi\in\mathfrak{g}$ durch
\begin{align}
	\Phi_{X_\xi}(t;g)
	&= \exp(t\xi)\,g
\end{align}
gegeben.

\subsubsection{Eigenschaften der Exponentialabbildung}

Wir notieren in diesem Abschnitt einige anaytischen und algebraischen Eigenschaften der Exponentialabbildung.

\paragraph{Einparametergruppe.}
Die Abbildung $t\mapsto\exp(t\xi)$ für festes $\xi\in\mathfrak{g}$ ist ein Gruppenhomomorphismus $(\mathbb{R},+)\rightarrow G$.

\paragraph{Ableitung der Einparametergruppe.}
Die Ableitung von $t\mapsto\exp(t\xi)$ lautet
\begin{align}
	\frac{\dt}{\dt t}\exp(t\xi)
	&= 
	\frac{\dt}{\dt t}
	\Phi_{X_\xi}(t;e)
	= X_\xi|_{\exp(t\xi)}
	= \Tt R_{\exp(t\xi)}|_e\cdot\xi
	\in \Tt_{\exp(t\xi)}G
	\ .
\end{align}
In der Parallelisierung entspricht dieses Ergebnis $\xi\in\mathfrak{g}$.

\paragraph{Ableitung im Ursprung.}
Für die Exponentialabbildung erhalten wir für $\xi\in\mathfrak{g}$ und $\upsilon\in\Tt_\xi\mathfrak{g}=\mathfrak{g}$
\begin{align}
\label{eq:dExp0}
	\Dt\exp|_\xi \cdot\upsilon
	&=
	\Tt R_{\exp(\xi)}^{-1}|_{\exp(\xi)}
	\cdot
	\frac{\dt}{\dt t}
		\exp(\xi + t\upsilon)
	\Big|_{t=0}
	\ ,
\end{align}
woraus für den Ursprung~$\xi=0$
\begin{align*}
	\Dt\exp|_0 \cdot\upsilon
	&=
	\frac{\dt}{\dt t}
		\exp(t\upsilon)
	\Big|_{t=0}
	= X_\upsilon|_e = \upsilon
	\ ,
\end{align*}
also
\begin{empheq}[box=\resultbox]{align}
	\Dt\exp|_0
	&=
	\mathrm{id}
\end{empheq}
folgt.

\paragraph{Normalkoordinaten.}
Da die Ableitung von $\exp$ im Ursprung gleich der Identität ist,
existiert ein Gebiet $U\subset\mathfrak{g}$, auf dem $\exp$ ein Diffeomorphismus auf $V=\exp U$ ist.
Die Umkehrabbildung
\begin{empheq}[box=\definitionbox]{align}
	\ln=\exp^{-1}
\end{empheq}
auf $V$ ist damit also eine lokale Karte von $G$ um das neutrale Element.
Über die (Rechts-)Multiplikation gewinnt man entsprechend eine lokale Karte um jedes $g\in G$.

Sei $\xi_1,\ldots,\xi_n$ eine Basis von $\mathfrak{g}$.
Indem wir Vektoren $\xi\in\mathfrak{g}$ bezüglich dieser Basis parametrisieren, $\xi=\sum_{i=1}^n x^i\xi_i$ mit $x^i\in\mathbb{R}$,
gewinnen wir lokale Koordinaten für $G$.
Diese nennen wir \mydef{(liesche) Normalkoordinaten}.

\subsubsection{Exponentialabbildung und adjungierte Darstellung}

Das Diagramm
\begin{center}
\begin{tikzcd}
\mathfrak{g} \arrow{d}{\exp} \arrow{r}{\ad}
&\mathfrak{gl}(\mathfrak{g}) \arrow{d}{\exp}\\
G \arrow{r}{\Ad}
& \mathrm{GL}(\mathfrak{g})
\end{tikzcd}
\end{center}
illustriert zwei Wege für Abbildungen $\mathfrak{g}\rightarrow\mathrm{GL}(\mathrm{g})$.
Dabei ist $\exp:\mathfrak{gl}(\mathfrak{g})\rightarrow\mathrm{GL}(\mathfrak{g})$ die ``gewöhnliche'' Operatorexponentialfunktion.

Dass sie übereinstimmen, dass also
\begin{empheq}[box=\resultbox]{align}
	\Ad\circ\exp &= \exp\circ\ad
\end{empheq}
ist, soll hier gezeigt werden.

Sei dazu $\xi\in\mathfrak{\xi}$ beliebig.
Wir betrachten zwei Kurven in $\mathrm{GL}(\mathfrak{g})$,
\begin{align*}
	\Gamma_1: t&\mapsto \exp(t\,\ad(\xi))
	\ ,\\
	\Gamma_2: t&\mapsto \Ad(\exp(t\xi))
	\ .
\end{align*}
$\Gamma_1$ ist nach Definition eine Einparametergruppe.
Dass dies auch $\Gamma_2$ ist, folgt aus
\begin{align*}
	\Gamma_2(t+s)
	&= \Ad(\exp((t+s)\xi))
	= \Ad(\exp(t\xi)\exp(s\xi))
	\nonumber\\
	&= \Ad(\exp(t\xi)\cdot\Ad(\exp(s\xi)
	\ .
\end{align*}
Es ist außerdem
\begin{align*}
	\dot{\Gamma}_1(0)
	&= \ad(\xi)
\end{align*}
und
\begin{align*}
	\dot{\Gamma}_2(0)
	&= \Dt\,\Ad|_e\cdot \xi = \ad(\xi) = \dot{\Gamma}_1(0)
	\ .
\end{align*}
Aufgrund der Einparametergruppeneigenschaft gilt für $i=1,2$
\begin{align*}
	\dot{\Gamma}_i(t)
	&=
	\frac{\dt}{\dt s}
	\Gamma_i(s+t)
	\Big|_{s=0}
	= \dot{\Gamma}_i(0)\Gamma_i(t)
	= \ad(\xi)\circ\Gamma_i(t)
\end{align*}
(wobei $\Gamma_i(t)$ und $\ad(\xi)$ vermöge der Einparametergruppe kommutieren).

Beide Kurven erfüllen also dieselbe gewöhnliche Differenzialgleichung ersten Grades und haben wegen $\Gamma_1(0)=\Gamma_2(0)=\mathrm{id}$ überdies den gleichen Anfangswert.
Damit stimmen $\Gamma_1$ und $\Gamma_2$ überein, 
so dass die Gleichheit $\Gamma_1(1)=\Gamma_2(1)$ die zu zeigende Aussage ergibt.

\subsubsection{Ableitung der Exponentialfunktion}

Ausgehend von Gl.~(\ref{eq:dExp0}) betrachten wir die Ableitung mit um den Faktor $s$ skalierten Tangentialvektoren,
\begin{align*}
	\Dt\exp|_{s\xi} \cdot s\upsilon
	&=
	\Tt R_{\exp(s\xi)}^{-1}|_{\exp(s\xi)}
	\cdot
	\frac{\dt}{\dt t}
		\exp(s\xi + st\upsilon)
	\Big|_{t=0}
	\nonumber\\
	&=
	\frac{\dt}{\dt t}
		\exp(s\xi + st\upsilon)
		\exp(-s\xi)
	\Big|_{t=0}
\end{align*}
Mit Hilfe der durch 
\begin{align*}
	h_{\xi,\upsilon}(t, s)
	&:= 
	\exp(s\xi + st\upsilon)\exp(-s\xi)
\end{align*}
definierten Abbildung $h_{\xi,\upsilon}:\mathbb{R}^2\rightarrow G$ schreiben wir
\begin{align*}
	\Dt\exp|_{s\xi} \cdot s\upsilon
	&= \dot{h}_{\xi, \upsilon}(0, s)
	\ ,
\end{align*}
wobei
\begin{align*}
	\dot{h}_{\xi, \upsilon}(t, s)
	&= \frac{\dt}{\dt t} h_{\xi, \upsilon}(t, s)
	\ .
\end{align*}
Wir notieren einige in der Folge wichtige Eigenschaften dieser Abbildung,
\begin{align*}
	h_{\xi,\upsilon}(0, s)
	&= e
	\ ,\\
	\dot{h}_{\xi, \upsilon}(0, s) 
	&\in\mathfrak{g}
	\ ,\\
	\dot{h}_{\xi, \upsilon}(0, 0)
	&= 0
	\ ,\\
	\dot{h}_{\xi, \upsilon}(0, 1)
	&= \Dt\exp|_{\xi} \cdot \upsilon
	\ .
\end{align*}
Die explizite Berechnung von $\dot{h}_{\xi, \upsilon}(0, s)$ erfolgt über die Ableitung nach $s$, beginnend mit der Produktregel (für $\Tt\,\mathrm{mult}$),
\begin{align*}
	\frac{\partial}{\partial s}
	\dot{h}_{\xi, \upsilon}(0, s)	
	&=
	\frac{\dt}{\dt t}	
	\Big(
		\Tt R_{\exp(-s\xi)}
		|_{\exp(s\xi + st\upsilon)}
		\cdot	
		\frac{\partial}{\partial s}\exp(s\xi + st\upsilon)
	\nonumber\\
	&\qquad\qquad+
		\Tt L_{\exp(s\xi + st\upsilon)}
		|_{\exp(-s\xi)}		
		\cdot			
		\frac{\partial}{\partial s}
		\exp(-s\xi)
	\Big)	
	\Big|_{t=0}
	\nonumber\\	
	&=
	\frac{\dt}{\dt t}	
	\Big(
		\Tt R_{\exp(-s\xi)}
		|_{\exp(s\xi + st\upsilon)}
		\cdot	
		\Tt R_{\exp(s\xi + st\upsilon)}
		|_e
		\cdot
		(\xi + t\upsilon)
	\nonumber\\
	&\qquad\qquad-
		\Tt L_{\exp(s\xi + st\upsilon)}
		|_{\exp(-s\xi)}	
		\cdot
		\Tt R_{\exp(-s\xi)}
		|_e
		\cdot
		\xi
	\Big)	
	\Big|_{t=0}
	\nonumber\\	
	&=
	\frac{\dt}{\dt t}	
	\Big(
		\Tt R_{h_{\xi, \upsilon}(t, s)}
		|_e
		\cdot
		(\xi + t\upsilon)
	\nonumber\\
	&\qquad\qquad-
		\Tt L_{h_{\xi, \upsilon}(t, s)}
		|_{e}	
		\cdot
		\Ad(\exp(-s\xi))
		\cdot
		\xi
	\Big)	
	\Big|_{t=0}
	\ .
\end{align*}
Mit
\begin{align*}
	\frac{\dt}{\dt t}	
	\Big(
		\Tt R_{h_{\xi, \upsilon}(t, s)}
		|_e
		\cdot
		t\upsilon
	\Big)
	\Big|_{t=0}
	&=
	\Tt R_{h_{\xi, \upsilon}(0, s)}
	|_e
	\cdot
	\upsilon
	= \upsilon
\end{align*}
und
\begin{align*}
	\Ad(\exp(-s\xi))\cdot\xi
	&= \exp(\ad(-s\xi))\cdot\xi
	= \xi
\end{align*}
erhalten wir 
\begin{align*}
	\frac{\partial}{\partial s}
	\dot{h}_{\xi, \upsilon}(0, s)	
	&=
	\upsilon +
	\frac{\dt}{\dt t}		
	\Big(
		\Tt R_{h_{\xi, \upsilon}(t, s)}
		|_e -
		\Tt L_{h_{\xi, \upsilon}(t, s)}
		|_{e}
	\Big)
	\Big|_{t=0}
	\cdot\xi
	\ .
\end{align*}
Mit
\begin{align*}
	\big(
	\Tt R_{h_{\xi, \upsilon}(t, s)}
	|_e -
	\Tt L_{h_{\xi, \upsilon}(t, s)}
	|_{e}
	\big)\big|_{t=0}
	&= 0
\end{align*}
sowie $\Tt L_{h_{\xi, \upsilon}(0, s)^{-1}}|_{h_{\xi, \upsilon}(0, s)}=\mathrm{id}$ (was trivial ist) folgt
\begin{align*}
	&\frac{\dt}{\dt t}	
	\Big(
		\Tt R_{h_{\xi, \upsilon}(t, s)}
		|_e -
		\Tt L_{h_{\xi, \upsilon}(t, s)}
		|_{e}	
	\Big)	
	\Big|_{t=0}
	\nonumber\\
	&=
	\frac{\dt}{\dt t}
	\big(
		\Tt L_{h_{\xi, \upsilon}(t, s)^{-1}}
		|_{h_{\xi, \upsilon}(t, s)} \cdot
		\Tt R_{h_{\xi, \upsilon}(t, s)}
		|_e -
		\mathrm{id}	
	\big)
	\Big|_{t=0}
	\nonumber\\
	&=
	\frac{\dt}{\dt t}
		\Ad(h_{\xi, \upsilon}(t, s)^{-1})
	\Big|_{t=0}
	= \ad(-\dot{h}_{\xi, \upsilon}(0, s))
	\ .
\end{align*}
Zusammengefasst erhalten wir
\begin{align*}
	\frac{\partial}{\partial s}
	\dot{h}_{\xi, \upsilon}(0, s)
	&=
	\ad(-\dot{h}_{\xi, \upsilon}(0, s))\cdot\xi + \upsilon
	= \ad(\xi)\cdot\dot{h}_{\xi, \upsilon}(0, s) + \upsilon
	\ .
\end{align*}
Dies ist eine gewöhnliche Diffenzialgleichung für $\dot{h}_{\xi, \upsilon}(0, s)$ auf dem endlichdimensionalen Vektorraum $\mathfrak{g}$.
Mit dem Anfangswert $\dot{h}_{\xi, \upsilon}(0, 0)=0$ finden wir
\begin{align*}
	\dot{h}_{\xi, \upsilon}(0, s)
	&=
	\int_0^s
	\exp((s-s')\,\ad(\xi))\cdot{\upsilon}
	\ \dt s'
\end{align*}
und mit $\dot{h}_{\xi, \upsilon}(0, 1)=\Dt\exp|_{\xi}\cdot\upsilon$ schließlich
\begin{empheq}[box=\resultbox]{align}
	\Dt\exp|_{\xi}
	&= E(\ad(\xi))
	\ ,
\end{empheq}
wobei
\begin{align}
	E(x) &:=
	\int_0^1
	\mathrm{e}^{(1-s')x}\ \dt s'
	= \frac{\mathrm{e}^x-1}{x}
	= 1 + \frac{1}{2}x + \frac{1}{6}x^2 + \ldots
	\ .
\end{align}

\section{Zusammenhänge auf Lie-Gruppen}

\subsection{Flacher Zusammenhang}

\subsection{Lie-Gruppen als riemannsche Mannigfaltigkeiten}

\subsection{Liesche und riemannsche Exponentialabbildungen}

\section{Anwendung in der Navigation}




\begin{thebibliography}{9}

\bibitem{Gallot2005}
Sylvestre Gallot, Dominique Hulin, Jacques Lafontaine:
\textit{Riemannian Geometry}.
3rd Edition, Universitext, Springer, Berlin/Heidelberg, 2005.

\bibitem{StorchWiebe2005}
Uwe Storch, Hartmut Wiebe:
\textit{Lehrbuch der Mathematik, Band 4: Analysis auf Mannigfaltigkeiten -- Funktionentheorie -- Funktionalanalysis}.
Spektrum Akademischer Verlag, Heidelberg, 2005.

\end{thebibliography}



\end{document}